\section{Lazy \texttt{let} and Dynamic Scope}
\label{sec:dynscope}

\code{lecture5b.rkt} and \code{lecture5c.rkt} make the environment interpreter \emph{lazy}, as \Cref{sec:lazy} did for substitution: \code{let} does not evaluate \code{e1}, but stores it and evaluates it only when \code{x} is used. The language is still that of \code{lecture5a.rkt}, numbers and \code{let}. The first attempt gets scope wrong; the second fixes it.

\subsection{Attempt 1: Store the Expression}

A binding now holds an unevaluated expression, and \code{lookup} returns it:
\begin{lstlisting}
(define-type Binding
  [bind (name : symbol) (defn : Expr)])

(define (lookup (x : symbol) (env : Env)) : Expr
  ...)   ; as in lecture5a, returning (bind-defn (first env))

(define (eval-env1 (env : Env) (e : Expr)) : number
  (type-case Expr e
    [numC (n) n]
    [plusC (e1 e2) (+ (eval-env1 env e1) (eval-env1 env e2))]
    [timesC (e1 e2) (* (eval-env1 env e1) (eval-env1 env e2))]
    [idC (x) (eval-env1 env (lookup x env))]
    [letC (x e1 e2) (eval-env1 (extend-env (bind x e1) env) e2)]))

(define (eval1 (e : Expr)) : number
  (eval-env1 empty-env e))
\end{lstlisting}
\code{letC} stores \code{e1} as is. \code{idC} looks up the stored expression and evaluates it, in \code{env}, the environment \emph{at the use}. The first examples work (``So far so good'', says the file):
\begin{lstlisting}[style=repl]
> (eval1 (parse '(+ 27 (let x 5 (+ x x)))))
37
> (eval1 (parse '(let x 5 (let y 6 (+ x y)))))
11
> (eval1 (parse '(let x 5 (let x 6 (+ x x)))))
12
\end{lstlisting}
Then the lecture tries a program where the stored expression mentions a name that is rebound before the use:
\begin{lstlisting}[style=repl]
(let x 1 (let y x (let x 2 (+ y y))))   env = []
  store x=1                             env = [x=1]
  store y=x            (the expression x, not a value)
                                        env = [y=x, x=1]
  store x=2                             env = [x=2, y=x, x=1]
  (+ y y)
    y ---> its expression is x; evaluate x in the CURRENT env
         ---> lookup x finds x=2 ---> 2
---> 4
\end{lstlisting}
``We get 4, but under our substitution semantics this would give us 2'': when \code{y} was defined, the only \code{x} in scope was \code{x=1}, so \code{y} means 1 (\Cref{fig:dyn}). There is a quicker way to see that 4 is wrong. Rename the inner \code{x} to \code{z}; the program is $\alpha$-equivalent (\Cref{sec:substrev}), so the answer must not change:
\begin{lstlisting}[style=repl]
> (eval1 (parse '(let x 1 (let y x (let x 2 (+ y y))))))
4
> (eval1 (parse '(let x 1 (let y x (let z 2 (+ y y))))))
2
\end{lstlisting}
It does change, so \code{eval-env1} is wrong. The file's last example is the counterpart of capture in \Cref{sec:lazy}:
\begin{lstlisting}[style=repl]
> (eval1 (parse '(let x y (let y 2 (+ x x)))))
4
\end{lstlisting}
The \code{y} in \code{(let x y ...)} is free, so this should be an error, but by the time \code{x} is used, the environment happens to contain \code{y=2}.

\begin{figure}[htbp]
\centering
\begin{tikzpicture}[tok/.style={inner sep=0pt, font=\ttfamily\normalsize, anchor=base west}]
  \node[tok] at (0,0) {(let x 1 (let y x (let x 2 (+ y y))))};
  \node[tok] (ox) at (5\cw,0) {\phantom{x}};
  \node[tok] (by) at (14\cw,0) {\phantom{x}};
  \node[tok] (ux) at (16\cw,0) {\phantom{x}};
  \node[tok] (ix) at (23\cw,0) {\phantom{x}};
  \node[tok] (u1) at (30\cw,0) {\phantom{x}};
  \node[tok] (u2) at (32\cw,0) {\phantom{x}};
  \draw[bind] (u1.south) to[out=-110, in=-70] (by.south);
  \draw[bind] (u2.south) to[out=-110, in=-70] (by.south);
  \draw[bind] (ux.south) to[out=-120, in=-60] (ox.south);
  \draw[bad] (ux.north) to[out=60, in=120] (ix.north);
  \node[lbl, text=accent] at (8\cw,-1.45) {lexical: the \code{x} in scope at \code{y}'s \code{let}};
  \node[lbl, text=sigc] at (19.5\cw,0.75) {dynamic: the newest \code{x} at the use};
\end{tikzpicture}
\caption{What the \code{x} stored in \code{y} refers to. Lexically (dashed blue) it is the outer \code{x}, giving 2; \code{eval-env1} evaluates it after the inner \code{let x 2} and finds that one instead (red), giving 4.}
\label{fig:dyn}
\end{figure}

\subsection{Dynamic vs.\ Lexical Scope}

\begin{definition}[Lexical and dynamic scope]
Under \emph{lexical} (static) scope, the binding a name refers to is determined by the program text: the nearest enclosing binder of that name, which can be found without running the program. Under \emph{dynamic} scope, it is determined by the execution history: the most recent binding of that name that is still active when the name is evaluated.
\end{definition}
The file adds that this is ``slightly informal, we could give a more precise definition''; the binding arrows of \Cref{fig:dyn} are the informal picture.

Substitution always gives lexical scope: it replaces names inside the binder's body at the moment the binder is processed, so a name can only ever be replaced by its own binder. \code{eval-env1} gives dynamic scope, because it evaluates the stored expression in whatever environment exists when the name is used, and by then other bindings may have been added.

\begin{pitfall}[Dynamic scope breaks $\alpha$-equivalence]
Under dynamic scope, renaming a local variable can change what \emph{other} code means, here the definition of \code{y}. A function's meaning depends on every caller's local names. This is why dynamic scope is considered a design mistake for ordinary variables, and why the test ``rename a bound variable and compare'' is such a useful bug detector.
\end{pitfall}

\subsection{Attempt 2: Store the Environment Too}

\code{lecture5c.rkt} ``corrects the issue by saving the environment \emph{at the time of the let} inside of a binding'', and evaluates the stored expression in that environment:
\begin{lstlisting}
(define-type Binding
  [bind (name : symbol) (env : Env) (defn : Expr)])

(define (lookup (x : symbol) (env : Env)) : Binding
  ...)   ; returns the whole binding, (first env)

(define (eval-env2 (env : Env) (e : Expr)) : number
  (type-case Expr e
    [numC (n) n]
    [plusC (e1 e2) (+ (eval-env2 env e1) (eval-env2 env e2))]
    [timesC (e1 e2) (* (eval-env2 env e1) (eval-env2 env e2))]
    [idC (x) (let ([res (lookup x env)])
               (eval-env2 (bind-env res) (bind-defn res)))]
    [letC (x e1 e2) (eval-env2 (extend-env (bind x env e1) env) e2)]))

(define (eval2 (e : Expr)) : number
  (eval-env2 empty-env e))
\end{lstlisting}
Two lines changed. \code{letC} stores \code{env}, the environment of the \code{let}, alongside \code{e1}. \code{idC} evaluates the stored expression in the stored environment, not the current one. A binding refers to an earlier environment, so the types are mutually recursive: a \tyname{Binding} contains an \tyname{Env}, which is a list of \tyname{Binding}s.

\begin{lstlisting}[style=repl]
(let x 1 (let y x (let x 2 (+ y y))))   env = []
  store x=(1 in [])                     env = [x]
  store y=(x in [x=1])                  env = [y, x]
  store x=(2 in [y, x=1])               env = [x, y, x]
  (+ y y)
    y ---> evaluate x in y's saved env [x=1] ---> 1
---> 2
\end{lstlisting}
``Now these examples have behaviors like what our substitution semantics produced'':
\begin{lstlisting}[style=repl]
> (eval2 (parse '(let x 1 (let y x (let x 2 (+ y y))))))
2
> (eval2 (parse '(let x 1 (let y x (let z 2 (+ y y))))))
2
> (eval2 (parse '(let x y (let y 2 (+ x x)))))
lookup: No binding found
\end{lstlisting}
In the last one, \code{x}'s expression \code{y} is evaluated in \code{x}'s saved environment, which is empty. The interpreter is still lazy: \code{(let x (+ y y) 10)} is 10, since \code{x} is never used.

\begin{keypoint}[The fix, in one sentence]
Run stored code in the environment where it was \emph{written}, not where it is \emph{used}. The pair (expression, environment) is the general tool for this; \Cref{sec:closures} needs it again for functions.
\end{keypoint}
